%% Supplementary Material — separate IEEEtran document (not counted in the
%% main paper's page limit). Body generated by build-supp.sh -> supp-body.tex.
%% Double-anonymous: no author identity, affiliations, funding, or links.
\documentclass[journal]{IEEEtran}

\usepackage{fontspec}
\setmainfont{TeX Gyre Termes}

\usepackage{amsmath,amssymb}
\usepackage{mathrsfs}
\usepackage[numbers,sort&compress]{natbib}
\usepackage{graphicx}
\usepackage{booktabs}
\usepackage{array}
\usepackage{multirow}
\usepackage{longtable}
\usepackage{calc}
\providecommand{\real}[1]{#1}
\usepackage{url}
\usepackage{xcolor}
\usepackage[colorlinks=true,linkcolor=blue,citecolor=blue,urlcolor=blue]{hyperref}

\providecommand{\tightlist}{\setlength{\itemsep}{0pt}\setlength{\parskip}{0pt}}
\providecommand{\passthrough}[1]{#1}
\providecommand{\pandocbounded}[1]{#1}

\setcounter{secnumdepth}{0}
\setkeys{Gin}{width=\linewidth,keepaspectratio}

\begin{document}

\title{Supplementary Material for\\
``Understanding Q-Learning and Deep Q-Learning in 2026:\\
A Methodological and Empirical Survey''}

\author{Anonymous Submission to IEEE Transactions on Artificial Intelligence%
\thanks{Double-anonymous review copy. This document supplements the main
manuscript and is not counted toward its page limit.}}

\markboth{IEEE Transactions on Artificial Intelligence --- Supplementary Material}%
{Understanding Q-Learning and Deep Q-Learning in 2026 --- Supplement}

\maketitle

\noindent
This document collects the auditable artifacts and extended material
referenced by the main paper: the systematic search log (supporting
§III), the full per-method method-type index (§III, §IV, §VIII), the
complete per-game Atari benchmark tables (§V), the full repository
coverage matrix (§VII), notation and selected derivations (§IV.I), and
the full configuration of the tabular experiment (§VI). Section and
table references of the form ``§X'' point to the main manuscript.

% ===================== BODY (generated by build-supp.sh) =====================
\input{supp-body.tex}

% ============================ BIBLIOGRAPHY ==================================
\bibliographystyle{IEEEtran}
\bibliography{refs}

\end{document}
